\subsection{本章小结}\label{subsec:logic-expansion-chain-summary}

本章至此已经建立了一套完整而独立的抽象逻辑系统。它并不是以公理表为起点，也不是以某一类具体对象为起点，而是以一条结构扩张链为起点：每一层只增加新的结构槽位，只扩大对象层的表达与可断言能力，而不回头改写较低层已经成立的 forcing 语义。正因为如此，本章的意义不在于提出又一套并列理论，而在于为全文提供一个统一的上位逻辑骨架，使后续所有具体模型都可以被解释为这一骨架的实现。

从最紧的角度看，本章完成了五次关键的逻辑提升。

第一，本章把对象层事实从“无索引真值”改写成了“状态上的可断言性”。一切对象层命题都不再以抽象全局世界为背景，而统一写成
\[
p\Vdash \varphi.
\]
因此，事实不再是预先给定的，而是相对于某个状态所稳定成立的内容。状态越细，forcing 内容越多；状态越粗，未定部分越大。这一点构成了整章的语义核心，也是后面所有高层结构共同依赖的基线。

第二，本章把对象的进入方式从“直接有值”改写成“先被可接受引用接纳，再进入评价域”。于是，对象层不再直接谈论“某物的值”，而先谈论“某个引用是否已进入当前状态的语义域”。这意味着，指称先于评价，对象先于取值。没有可接受引用，就没有对象层的值命题；这使后续所有地址、前缀、标签、局部坐标与证书系统，都能够被统一解释为“引用系统”的不同实例，而不必各自再发明一套基础语义。

第三，本章把“对象存在”从单点存在改写成了局部片段、相容族与整体粘接的链式形成过程。一个对象之所以成为对象，不再是因为某处有一个立即可取的值，而是因为它首先在某些局部片段上形成局部对象，接着这些局部对象彼此相容，最后它们能够被粘接成整体对象。于是，本章第一次在逻辑层上区分了：局部存在、相容局部存在与整体存在。后续任何关于局部截面、局部证书与整体读出的讨论，都只能被看作这一局部—整体逻辑的具体化，而不再是与本章并列的结构。

第四，本章把 \(\mathsf{NULL}\) 从“特殊值”彻底改写成了“结构性失败标签”。对象处于 \(\mathsf{NULL}\) 状态，不是因为它取到了某个空值，而是因为它在局部形成链上发生了失败：要么局部片段尚未形成，要么局部片段形成了但彼此不相容，要么局部片段相容了但仍未能粘接成整体对象。于是，\(\mathsf{NULL}\) 被分解为局部空缺、相容空缺与粘接空缺三类模式。更重要的是，这一分解并不是一条无条件的经典析取，而是只有在状态已经足以区分中间层性质时，才会被真正展开。这一构造性立场保证了：失败模式本身也是需要信息支撑才能被判定的对象层事件。

第五，本章最终把状态 forcing 提升为一个动态的、可组合的、观察者时空化的系统。状态不再是静止的，而沿更新链不断细化；时间不再是外在背景，而是状态推进的顺序投影；多轴结构说明细化不必沿单一方向发生，而可由多个独立方向共同生成；观察者索引把状态空间分纤维化为不同局部经验链；事件与区域系统则把这些局部状态进一步嵌入到因果有序的时空背景之中。于是，forcing 的最终形式不再只是某个匿名状态上的成立，而是某观察者在某事件、某区域内对某命题的局部可断言性。

因此，若把本章的全部结构压缩为若干最核心的公式，它们可以写成：
\begin{itemize}
  \item 对象层事实：
  \[
  p\Vdash \varphi.
  \]
  \item 可接受引用：
  \[
  \mathsf{Adm}_s(r).
  \]
  \item 整体局部截面：
  \[
  \mathsf{Sec}_s(r).
  \]
  \item 结构性空缺：
  \[
  \mathsf{Null}_s(r)
  :=
  \mathsf{Adm}_s(r)\land \neg \mathsf{Sec}_s(r).
  \]
  \item 状态推进：
  \[
  q\in \mathrm{Upd}_\alpha(p)
  \qquad\text{或更抽象地}\qquad
  p\rightsquigarrow_K q.
  \]
  \item 观察者纤维：
  \[
  p\in P_i.
  \]
  \item 事件与区域定位：
  \[
  \lambda(p)=x,
  \qquad
  \varrho(p)\subseteq U.
  \]
  \item 共享 forcing：
  \[
  c\Vdash^{\forall}_{K,U}\varphi.
  \]
\end{itemize}

这几条公式实际上已经概括了整章的全部骨架。它们表明：本章最终建立的，不是一套额外的对象理论，而是一套统一的语义接口。后续具体章节若要进入全文，就必须回答同一个问题：它们各自究竟实现了这些抽象结构中的哪一部分。

于是，本章与后文的关系也可以被简洁地重述如下。任何具体的地址系统，都是本章引用逻辑的实例；任何具体的前缀站点与局部截面，都是本章局部性逻辑的实例；任何具体的 \(\mathsf{NULL}\) 三分解与失败证书，都是本章失败模式逻辑的实例；任何具体的更新链、时间推进、多轴残差结构，都是本章动态与多轴逻辑的实例；任何具体的观察者背景、事件定位与区域传播，则都是本章观察者时空逻辑的实例。后文不再需要另立逻辑起点，而只需要说明各自如何填充本章已经建立的结构槽位。

在这一意义上，本章的最终结论可以压缩为：\emph{本文不以公理表定义逻辑，而以结构扩张链定义逻辑；不以全局真值定义事实，而以状态 forcing 定义事实；不以单一对象系统组织全文，而以一个可实例化的观察者时空逻辑系统作为全文的统一上位框架。}

这句话既是本章的结语，也是全文后续模型化展开的总前提。接下来的具体章节，无论讨论递归地址化、局部截面、\(\mathsf{NULL}\) 三分解、多轴前缀还是观察结构，都应被理解为本章抽象逻辑系统的实现，而不是与本章并列的第二套基础。
